\chapter{The Interview}

\lettrine{R}{aoul} made one step towards the girl who thus called him.

\zz
<But my horse, madame?> said he.

\zz
<Oh! you are terribly embarrassed! Go yonder way—there is a shed in the outer court: fasten your horse, and return quickly!>

\zz
<I obey, madame.>

Raoul was not four minutes in performing what he had been directed to do; he returned to the little door, where, in the gloom, he found his mysterious conductress waiting for him, on the first steps of a winding staircase.

<Are you brave enough to follow me, monsieur knight errant?> asked the girl, laughing at the momentary hesitation Raoul had manifested.

The latter replied by springing up the dark staircase after her. They thus climbed up three stories, he behind her, touching with his hands, when he felt for the banister, a silk dress which rubbed against each side of the staircase. At every false step made by Raoul, his conductress cried, <Hush!> and held out to him a soft perfumed hand.

<One would mount thus to the belfry of the castle without being conscious of fatigue,> said Raoul.

<All of which means, monsieur, that you are very much perplexed, very tired, and very uneasy. But be of good cheer, monsieur; here we are, at our destination.>

The girl threw open a door, which immediately, without any transition, filled with a flood of light the landing of the staircase, at the top of which Raoul appeared, holding fast by the balustrade.

The girl continued to walk on—he followed her; she entered a chamber—he did the same.

As soon as he was fairly in the net he heard a loud cry, and, turning round, saw at two paces from him, with her hands clasped and her eyes closed, that beautiful fair girl with blue eyes and white shoulders, who, recognizing him, called him Raoul.

He saw her, and divined at once so much love and so much joy in the expression of her countenance, the he sank on his knees in the middle of the chamber, murmuring, on his part, the name of Louise.

<Ah! Montalais!—Montalais!> she sighed, <it is very wicked to deceive me so.>

<Who, I\@? I have deceived you?>

<Yes; you told me you would go down to inquire the news, and you have brought up monsieur!>

<Well, I was obliged to do so—how else could he have received the letter you wrote him?> And she pointed with her finger to the letter which was still upon the table.

Raoul made a step to take it; Louise, more rapid, although she had sprung forward with a sufficiently remarkable physical hesitation, reached out her hand to stop him. Raoul came in contact with that trembling hand, took it within his own, and carried it so respectfully to his lips, that he might have been said to have deposited a sigh upon it rather than a kiss.

In the meantime, Mademoiselle de Montalais had taken the letter, folded it carefully, as women do, in three folds, and slipped it into her bosom.

<Don't be afraid, Louise,> said she; <monsieur will no more venture to take it hence than the defunct king \regnal{Louis XIII} ventured to take billets from the corsage of Mademoiselle de Hautefort.>

Raoul blushed at seeing the smile of the two girls; and he did not remark that the hand of Louise remained in his.

<There!> said Montalais, <you have pardoned me, Louise, for having brought monsieur to you; and you, monsieur, bear me no malice for having followed me to see mademoiselle. Now, then, peace being made, let us chat like old friends. Present me, Louise, to M.~de Bragelonne.>

<Monsieur le Vicomte,> said Louise, with her quiet grace and ingenuous smile, <I have the honour to present to you Mademoiselle Aure de Montalais, maid of honour to her royal highness Madame, and moreover my friend—my excellent friend.>

Raoul bowed ceremoniously.

<And me, Louise,> said he—<will you not present me also to mademoiselle?>

<Oh, she knows you—she knows all!>

This unguarded expression made Montalais laugh and Raoul sigh with happiness, for he interpreted it thus: <She knows all our love.>

<The ceremonies being over, Monsieur le Vicomte,> said Montalais, <take a chair, and tell us quickly the news you bring flying thus.>

<Mademoiselle, it is no longer a secret; the king, on his way to Poitiers, will stop at Blois, to visit his royal highness.>

<The king here!> exclaimed Montalais, clapping her hands. <What! are we going to see the court? Only think, Louise—the real court from Paris! Oh, good heavens! But when will this happen, monsieur?>

<Perhaps this evening, mademoiselle; at latest, to-morrow.>

Montalais lifted her shoulders in a sigh of vexation.

<No time to get ready! No time to prepare a single dress! We are as far behind the fashions as the Poles. We shall look like portraits from the time of \regnal{Henry IV}. Ah, monsieur! this is sad news you bring us!>

<But, mesdemoiselles, you will be still beautiful!>

<That's no news! Yes, we shall always be beautiful, because nature has made us passable; but we shall be ridiculous, because the fashion will have forgotten us. Alas! ridiculous! I shall be thought ridiculous—I!>

<And by whom?> said Louise, innocently.

<By whom? You are a strange girl, my dear. Is that a question to put to me? I mean everybody; I mean the courtiers, the nobles; I mean the king.>

<Pardon me, my good friend; but as here every one is accustomed to see us as we are\longdash>

<Granted; but that is about to change, and we shall be ridiculous, even for Blois; for close to us will be seen the fashions from Paris, and they will perceive that we are in the fashion of Blois! It is enough to make one despair!>

<Console yourself, mademoiselle.>

<Well, so let it be! After all, so much the worse for those who do not find me to their taste!> said Montalais, philosophically.

<They would be very difficult to please,> replied Raoul, faithful to his regular system of gallantry.

<Thank you, Monsieur le Vicomte. We were saying, then, that the king is coming to Blois?>

<With all the court.>

<Mesdemoiselles de Mancini, will they be with them?>

<No, certainly not.>

<But as the king, it is said, cannot do without Mademoiselle Mary?>

<Mademoiselle, the king must do without her. M.~le Cardinal will have it so. He has exiled his nieces to Brouage.>

<He!—the hypocrite!>

<Hush!> said Louise, pressing a finger on her friend's rosy lips.

<Bah! nobody can hear me. I say that old Mazarino Mazarini is a hypocrite, who burns impatiently to make his niece Queen of France.>

<That cannot be, mademoiselle, since M.~le Cardinal, on the contrary, had brought about the marriage of his majesty with the Infanta Maria Theresa.>

Montalais looked Raoul full in the face, and said, <And do you Parisians believe in these tales? Well! we are a little more knowing than you, at Blois.>

<Mademoiselle, if the king goes beyond Poitiers and sets out for Spain; if the articles of the marriage contract are agreed upon by Don Luis de Haro and his eminence, you must plainly perceive that it is not child's play.>

<All very fine! but the king is king, I suppose?>

<No doubt, mademoiselle; but the cardinal is the cardinal.>

<The king is not a man, then! And he does not love Mary Mancini?>

<He adores her.>

<Well, he will marry her then. We shall have war with Spain. M.~Mazarin will spend a few of the millions he has put away; our gentlemen will perform prodigies of valour in their encounters with the proud Castilians, and many of them will return crowned with laurels, to be recrowned by us with myrtles. Now, that is my view of politics.>

<Montalais, you are wild!> said Louise, <and every exaggeration attracts you as light does a moth.>

<Louise, you are so extremely reasonable, that you will never know how to love.>

<Oh!> said Louise, in a tone of tender reproach, <don't you see, Montalais? The queen-mother desires to marry her son to the Infanta; would you wish him to disobey his mother? Is it for a royal heart like his to set such a bad example? When parents forbid love, love must be banished.>

And Louise sighed: Raoul cast down his eyes, with an expression of constraint. Montalais, on her part, laughed aloud.

<Well, I have no parents!> said she.

<You are acquainted, without doubt, with the state of health of M.~le Comte de la Fère?> said Louise, after breathing that sigh which had revealed so many griefs in its eloquent utterance.

<No, mademoiselle,> replied Raoul, <I have not let paid my respects to my father; I was going to his house when Mademoiselle de Montalais so kindly stopped me. I hope the comte is well. You have heard nothing to the contrary, have you?>

<No, M.~Raoul—nothing, thank God!>

Here, for several instants, ensued a silence, during which two spirits, which followed the same idea, communicated perfectly, without even the assistance of a single glance.

<Oh, heavens!> exclaimed Montalais in a fright; <there is somebody coming up.>

<Who can it be?> said Louise, rising in great agitation.

<Mesdemoiselles, I inconvenience you very much. I have, without doubt, been very indiscreet,> stammered Raoul, very ill at ease.

<It is a heavy step,> said Louise.

<Ah! if it is only M.~Malicorne,> added Montalais, <do not disturb yourselves.>

Louise and Raoul looked at each other to inquire who M.~Malicorne could be.

<There is no occasion to mind him,> continued Montalais; <he is not jealous.>

<But, mademoiselle\longdash> said Raoul.

<Yes, I understand. Well, he is discreet as I am.>

<Good heavens!> cried Louise, who had applied her ear to the door, which had been left ajar; <it is my mother's step!>

<Madame de Saint-Rémy! Where shall I hide myself?> exclaimed Raoul, catching at the dress of Montalais, who looked quite bewildered.

<Yes,> said she; <yes, I know the clicking of those pattens! It is our excellent mother. M.~le Vicomte, what a pity it is the window looks upon a stone pavement, and that fifty paces below it.>

Raoul glanced at the balcony in despair. Louise seized his arm and held it tight.

<Oh, how silly I am!> said Montalais; <have I not the robe-of-ceremony closet? It looks as if it were made on purpose.>

It was quite time to act; Madame de Saint-Rémy was coming up at a quicker pace than usual. She gained the landing at the moment when Montalais, as in all scenes of surprises, shut the closet by leaning with her back against the door.

<Ah!> cried Madame de Saint-Rémy, <you are here, are you, Louise?>

<Yes, madame,> replied she, more pale than if she had committed a great crime.

<Well, well!>

<Pray be seated, madame,> said Montalais, offering her a chair, which she placed so that the back was towards the closet.

<Thank you, Mademoiselle Aure—thank you. Come, my child, be quick.>

<Where do you wish me to go, madame?>

<Why, home, to be sure; have you not to prepare your toilette?>

<What did you say?> cried Montalais, hastening to affect surprise, so fearful was she that Louise would in some way commit herself.

<You don't know the news, then?> said Madame de Saint-Rémy.

<What news, madame, is it possible for two girls to learn up in this dove-cote?>

<What! have you seen nobody?>

<Madame, you talk in enigmas, and you torment us at a slow fire!> cried Montalais, who, terrified at seeing Louise become paler and paler, did not know to what saint to put up her vows.

At length she caught an eloquent look of her companion's, one of those looks which would convey intelligence to a brick wall. Louise directed her attention to a hat—Raoul's unlucky hat, which was set out in all its feathery splendour upon the table.

Montalais sprang towards it, and, seizing it with her left hand, passed it behind her into the right, concealing it as she was speaking.

<Well,> said Madame de Saint-Rémy, <a courier has arrived, announcing the approach of the king. There, mesdemoiselles; there is something to make you put on your best looks.>

<Quick, quick!> cried Montalais. <Follow Madame your mother, Louise; and leave me to get ready my dress of ceremony.>

Louise arose; her mother took her by the hand, and led her out on to the landing.

<Come along,> said she; then adding in a low voice, <When I forbid you to come the apartment of Montalais, why do you do so?>

<Madame, she is my friend. Besides, I had but just come.>

<Did you see nobody concealed while you were there?>

<Madame!>

<I saw a man's hat, I tell you—the hat of that fellow, that good-for-nothing!>

<Madame!> repeated Louise.

<Of that do-nothing Malicorne! A maid of honour to have such company—fie! fie!> and their voices were lost in the depths of the narrow staircase.

Montalais had not missed a word of this conversation, which echo conveyed to her as if through a tunnel. She shrugged her shoulders on seeing Raoul, who had listened likewise, issue from the closet.

<Poor Montalais!> said she, <the victim of friendship! Poor Malicorne, the victim of love!>

She stopped on viewing the tragic-comic face of Raoul, who was vexed at having, in one day, surprised so many secrets.

<Oh, mademoiselle!> said he; <how can we repay your kindness?>

<Oh, we will balance accounts some day,> said she. <For the present, begone, M.~de Bragelonne, for Madame de Saint-Rémy is not over indulgent; and any indiscretion on her part might bring hither a domiciliary visit, which would be disagreeable to all parties.>

<But Louise—how shall I know\longdash>

<Begone! begone! King \regnal{Louis XI} knew very well what he was about when he invented the post.>

<Alas!> sighed Raoul.

<And am I not here—I, who am worth all the posts in the kingdom? Quick, I say, to horse! so that if Madame de Saint-Rémy should return for the purpose of preaching me a lesson on morality, she may not find you here.>

<She would tell my father, would she not?> murmured Raoul.

<And you would be scolded. Ah, vicomte, it is very plain you come from court; you are as timid as the king. \swear{Peste!} at Blois we contrive better than that, to do without papa's consent. Ask Malicorne else!>

And at these words the girl pushed Raoul out of the room by the shoulders. He glided swiftly down to the porch, regained his horse, mounted, and set off as if he had had Monsieur's guards at his heels. 